%======================================================================
\chapter{Simulation Results}
\label{ch:results}
%======================================================================

%----------------------------------------------------------------------
\section{Simulation Setup}

All experiments use the default configuration in \texttt{config.py}
unless stated otherwise.
Table~\ref{tab:sim-params} summarizes the key parameters.

\begin{table}[htbp]
\centering
\caption{Simulation parameters}
\label{tab:sim-params}
\renewcommand{\arraystretch}{1.4}
\begin{tabular}{llll}
\toprule
Parameter & Symbol & Value & Units \\
\midrule
Domain size          & $L_x \times L_y$       & $5 \times 5$   & cm \\
Grid resolution      & $N_x \times N_y$       & $50 \times 50$ & --- \\
Time step            & $\Delta t$             & $1$            & s \\
Simulation horizon   & $t_f$                  & $300$          & s \\
Tissue density       & $\rho$                 & $1050$         & \si{\kilogram\per\cubic\meter} \\
Specific heat        & $c$                    & $3600$         & \si{\joule\per\kilogram\per\kelvin} \\
Thermal conductivity & $k$                    & $0.51$         & \si{\watt\per\meter\per\kelvin} \\
Perfusion rate       & $\omega_b$             & $0.005$        & \si{\per\second} \\
Max power            & $P_{\max}$             & $50$           & W \\
Safety temperature   & $T_{\text{safe}}$      & $45$           & \si{\degreeCelsius} \\
Tumor radius         & $r_T$                  & $8$            & mm \\
SAR beam width       & $\sigma_{\text{SAR}}$  & $6$            & mm \\
\bottomrule
\end{tabular}
\end{table}

%----------------------------------------------------------------------
\section{Open-Loop Baseline}

Figure~\ref{fig:openloop-T} shows the temperature field at $t_f = 300$~s
under the bang-bang open-loop schedule.

% ── PLACEHOLDER ──────────────────────────────────────────────────────
% Replace with: \includegraphics[width=0.7\textwidth]{figures/T_final_openloop.png}
\begin{figure}[htbp]
    \centering
    \fbox{\parbox{0.65\textwidth}{\centering\vspace{2cm}
        \textit{[Insert: T\_final.png — temperature field at $t_f$]} \\[0.5em]
        Generated by: \texttt{python main.py --mode openloop}
    \vspace{2cm}}}
    \caption{Final temperature field under open-loop bang-bang schedule.
             Colormap: \texttt{hot}. Tumor boundary overlaid in white.}
    \label{fig:openloop-T}
\end{figure}

Figure~\ref{fig:openloop-damage} shows the Arrhenius damage field
$\Omega_d(\mathbf{r}, t_f)$ with the $\Omega_d = 1$ necrosis boundary.

\begin{figure}[htbp]
    \centering
    \fbox{\parbox{0.65\textwidth}{\centering\vspace{2cm}
        \textit{[Insert: damage\_final.png — Arrhenius damage at $t_f$]} \\[0.5em]
        Generated by: \texttt{python main.py --mode openloop}
    \vspace{2cm}}}
    \caption{Arrhenius damage field $\Omega_d$ at $t_f$.
             White contour: necrosis boundary ($\Omega_d = 1$).
             Dashed circle: tumor boundary.}
    \label{fig:openloop-damage}
\end{figure}

%----------------------------------------------------------------------
\section{Indirect (PMP) Solution}

\begin{figure}[htbp]
    \centering
    \fbox{\parbox{0.65\textwidth}{\centering\vspace{2cm}
        \textit{[Insert: power\_history.png — optimized power schedule]} \\[0.5em]
        Generated by: \texttt{python main.py --mode indirect}
    \vspace{2cm}}}
    \caption{Optimized power schedule $P^*(t)$ from the indirect PMP solver.}
    \label{fig:indirect-power}
\end{figure}

%----------------------------------------------------------------------
\section{Solver Comparison}

Table~\ref{tab:solver-comparison} compares all four solver modes on
cost, ablation completeness, and peak healthy-tissue temperature.

\begin{table}[htbp]
\centering
\caption{Solver comparison (fill in from simulation output)}
\label{tab:solver-comparison}
\renewcommand{\arraystretch}{1.4}
\begin{tabular}{lllll}
\toprule
Solver & $J_{\text{total}}$ & Ablation & $T_{\max}$ ($\Omega_H$) & Runtime \\
\midrule
Open-loop (bang-bang) & ---  & --- \% & --- \si{\degreeCelsius} & ---~s \\
Indirect (PMP)        & ---  & --- \% & --- \si{\degreeCelsius} & ---~s \\
Direct (SLSQP)        & ---  & --- \% & --- \si{\degreeCelsius} & ---~s \\
MPC                   & ---  & --- \% & --- \si{\degreeCelsius} & ---~s \\
\bottomrule
\end{tabular}
\end{table}

%----------------------------------------------------------------------
\section{Boundary Condition Comparison}

The effect of boundary conditions on the steady-state temperature profile
is examined by running the open-loop mode with each of the three BC types.

\begin{figure}[htbp]
    \centering
    \fbox{\parbox{0.65\textwidth}{\centering\vspace{2cm}
        \textit{[Insert: BC comparison figure]} \\[0.5em]
        Generated by: \texttt{--bc dirichlet}, \texttt{--bc neumann}, \texttt{--bc robin}
    \vspace{2cm}}}
    \caption{Final temperature field under Dirichlet, Neumann, and Robin BCs
             (open-loop mode, same power schedule).
             Robin parameters: $h_c = \SI{50}{\watt\per\square\meter\per\kelvin}$,
             $T_\infty = \SI{20}{\degreeCelsius}$.}
    \label{fig:bc-comparison}
\end{figure}

%----------------------------------------------------------------------
\section{Adjoint Gradient Verification}

The adjoint gradient $\nabla_P J$ is verified against a centered
finite-difference approximation at a random control perturbation:
\begin{equation}
    \left.\frac{\partial J}{\partial P_k}\right|_{\text{FD}}
    \approx \frac{J(\mathbf{u} + \epsilon\,\mathbf{e}_k)
                  - J(\mathbf{u} - \epsilon\,\mathbf{e}_k)}{2\epsilon},
    \qquad \epsilon = 10^{-3}
\end{equation}

% Fill in relative error table after running gradient check
\noindent
Maximum relative error between adjoint and FD gradients: \textit{[fill in]} \%.
